\documentclass[11pt,a4paper]{article}
\usepackage[margin=1in]{geometry}
\usepackage{hyperref}
\usepackage{xcolor}

\title{Disclaimer for Kerr-Gordon Vortex Model}
\author{Joshua Christopher Ryan}
\date{\today}

\begin{document}

\maketitle

\section{Disclaimer}

The term ``Kerr-Gordon vortex'' is used in this work purely as a \textbf{creative and analogical label} for the effective lensing mechanism presented in the model.

It is inspired by the intuitive behavior of a \textbf{kaleidoscope magnifying glass} — involving multiple reflections, symmetry, magnification, and complex light-bending paths — rather than claiming a full rigorous derivation from the exact Kerr metric combined with a Gordon scalar field in general relativity.

This is a phenomenological toy model designed to explore possible resolutions to the Hubble tension through lensing-inspired corrections to the effective Hubble parameter. It is not presented as a complete or formally derived solution within standard general relativity or modified gravity frameworks.

All credit for the development and implementation of this model belongs to \textbf{Joshua Christopher Ryan}.

\vspace{1cm}
For questions or further exploration, feel free to contact the author or refer to the accompanying Python starter pack.

\end{document}
